\section{The seven lower-moment comparison plans}\label{app:finite-details}

This appendix gives the complete plans underlying Proposition~\ref{fc:lowprofiles}.
All use the common comparison, curvature and allocation rules of the main text.

\begin{lemma}[One lower-moment profile]\label{fc:low40}
Under the retained actual-$26$ acquisition and the global predecessor
$C_{25}\le\CE$, the entire original profile
$p\in[.40,.41]$, $1\le\rho\le11/10$ has a strictly negative
conditional count residual. A uniform upper bound is
$-5610547/10^{12}$.
\end{lemma}
\begin{proof}
The original inward and noise intercepts are
$\Delta=.057809$ and $\Lambda=.301161$.
The necessary signed resource plane improves the actual-noise cap
to $L_*=.17805629$. The original contact inequalities exclude
$25$ of the $27$ count roots, leaving only $10$ and $11$.

Current-count comparisons cover count $10$ on $[-1/2,0]$ with
slope multiplier $1$, and count $11$ on its two half intervals
with multipliers $3/2$ and $1$. Their upper bounds are respectively
$-.000005610547$, $-.000934723659$, and $-.003337615615$.

For the positive half of count $10$, retain the seven noise
intervals and four allocation intervals above.
On the first noise interval join the current price to the
reflected-square price at $t=1/32$; on the second join at
$1/8$; on intervals $3$--$7$ join at $1/4$.
Each reflected-square price uses reference count $11$, width
$h=7/8$ and Corollary~\ref{ce:allocation-interval}.
Here $11>26p$, $p<1/2$, and the positive affine inward cap
decreases in allocation. The moment floor therefore increases
in allocation. All next-atom, legal-width and same-noise derivative
gates hold on their whole stated domains. In the first interval
$y_{\max}=3/32>0$ is allowed by the extended variance-derivative
bound in Section~\ref{pp:section}.

Every current prefix is at most $-.000025632850$;
every clipped suffix is at most $-.001597137326$.
Probability and phase chords include every closed boundary.
The two pieces in each noise/allocation region meet at the stated
phase, and the full noise union is $[0,L_*]$.
Together with the three half-interval comparisons and the contact
exclusions this covers the whole original profile. The least strict
bound is $-.000005610547$.
\end{proof}

This replaces a conditional count comparison. The original support,
floor and acquisition proofs remain its inputs; replaying the pure
physical acquisition formulas does not establish those inputs anew.

\begin{lemma}[An adjacent lower-moment profile]\label{fc:low39}
Under the retained actual-$26$ acquisition and $C_{25}\le\CE$,
the entire original profile $p\in[.39,.40]$, $1\le\rho\le11/10$
has conditional count residual at most $-98710877/10^{12}$.
\end{lemma}
\begin{proof}
The original intercepts are $\Delta=.05769$, $\Lambda=.30708$;
the same signed resource argument gives $\lambda\le .17375152$.
Contact excludes $24$ count roots. Only $9,10,11$ remain, and
the count-$9$ contact lower phase is greater than $1/4$.
Count $11$ has the two current-count half-interval comparisons
with slope multipliers $3/2,1$.

For the lower counts one reference law covers both sides of a
nearest-count boundary. Use source count $9$, reference count $10$
and width $7/8$. On noise intervals $1,2$ its auxiliary phase runs
through $[1/4,7/8]$, and on intervals $3$--$7$ through $[1/4,3/4]$.
The exact threshold identity
\[
 9+(1+t)=10+t
\]
therefore covers count $9$ and the lower part of count $10$ by
one continuous price in each noise group. Lemma~\ref{mc:staircase}
holds for every real phase, so an auxiliary phase greater than
$1/2$ is permitted. Here $10<26p$ throughout the probability
interval, and Corollary~\ref{ce:allocation-interval} selects the
lower allocation endpoint for both variance and moment cost.

The remaining count-$10$ negative phases are $[-1/8,0]$ on
the first two noise intervals and $[-1/4,0]$ on the others;
current-count prices cover them. Its positive half uses the
same $1/32,1/8,1/4$ handoffs as in Lemma~\ref{fc:low40}, now
with this profile's own resources. Reference count $11>26p$
selects the upper allocation endpoint for variance and the lower
one for moment cost. All legal-width, next-atom, noise-derivative
and continuous probability/phase gates are strict on their stated
domains. The whole closed unions retain every original count,
noise and allocation region, with largest upper bound
$-.000098710877$.
\end{proof}

\begin{lemma}[Reversed signed noise prices]\label{fc:low41}
Under the retained actual-$26$ acquisition and $C_{25}\le\CE$,
the entire original profile $p\in[.41,.42]$, $1\le\rho\le11/10$
has conditional count residual at most $-24705293/(5\cdot10^{11})$.
\end{lemma}
\begin{proof}
The original intercepts are $\Delta=.054236$, $\Lambda=.276951$.
Here $h_+-h_-<0$. The positive endpoint coefficients in
\eqref{loc:two-endpoint-noise} still give a valid finite resource
iteration and $\lambda\le L_*=.162619$, retaining the original
joint intercept $\Lambda$. For the actual moment floor,
\eqref{ce:signed-allocation-gate} is greater than $.473372164977$;
both negative terms from the increasing inward cap are included.
Thus $G$ increases with allocation on the whole domain.

For current-count prices, the necessary weaker plane
\[
 g(a_0+a_1)+\lambda h_-
 \le\mathfrak R+L_*(h_--h_+)_+
\]
supplies an allocation-independent inward cap. Its intercept is
smaller than the original $\Delta$, and its noise slope is larger
than $\Delta/\Lambda$, as required by the fixed-reference estimate.

Contact excludes $24$ count roots and leaves $10,11,12$.
Count $10$ starts above $-1/8$, and count $12$ ends below $-1/4$.
Current-count prefixes cover count $10$ through $1/32,1/8,1/4$
on the first, second and remaining noise intervals.
One source-$10$, reference-$11$ square then covers both its
positive suffix and the negative prefix of count $11$.
On noise intervals $1,2$ the width is $7/8$ and its auxiliary
phase ends at $7/8$. On intervals $3$--$6$ use width $9/8$
with phase $[1/4,7/8]$. On the last interval use width $7/8$
and phase $[1/4,3/4]$. The identity $10+(1+t)=11+t$
gives the original count-$11$ phases exactly.

Current-count prices complete count $11$ through zero and on its
positive half; count $12$ uses slope multiplier $3/2$ on
$[-1/2,-1/4]$. The original seven noise intervals and four
allocation intervals are unchanged. All width, next-atom,
allocation and noise gates, followed by the probability/phase
chords, give a strict bound on the full closed unions.
Their largest upper bound is $-.000049410586$.
\end{proof}

\begin{corollary}[A complete lower-moment strip]\label{fc:p41-strip}\label{fc:n27-low41}\label{fc:n28-low41}\label{fc:n29-low41}\label{fc:n30-low41}
At each integer $n=26,\ldots,30$, under the retained actual-$n$ acquisition
and $C_{n-1}\le\CE$, the original profile $p\in[.41,.42]$,
$1\le\rho\le11/10$ has the following conditional count bounds.
The original noise intercept is $\Lambda=.276951$ in all rows.
The five residual bounds are the corresponding rows of
Proposition~\ref{fc:lowprofiles}.

\end{corollary}
\begin{proof}
Use the same comparison rules with each integer's own masses,
variance and moment normalization. Lemma~\ref{fc:low41} gives $26$.
At $27$, contact excludes
$25$ of the $28$ roots and leaves $10,11,12$.
Reference count $11<27p$ uses lower allocation
endpoints for variance and cost, while $12>27p$ uses the two
allocation corners. The signed moment derivative gate is positive.
On the lowest noise interval the source-$10$ square
runs through auxiliary phase $15/16$, which is exactly count-$11$
phase $-1/16$. The current price on $[-1/16,1/32]$ and its
clipped-square continuation then close the middle count.

At $28$, the remaining counts are $11$ and $12$.
Their outer phase intervals are $[-17/64,1/2]$ and $[-1/2,15/32]$.
The same fixed current
references for count $11$ use the phase faces $-1/8$ and $0$,
followed by the square continuations. Count $12$ joins current
slopes $3/2$ and $1$ at phase $-1/32$. Reference count $12>28p$
has fixed variance direction throughout the probability interval.

At $29$, contact excludes $27$ roots and leaves $11,12,13$.
Reference count $12$ crosses $29p$ at $p=12/29$ inside the
original probability interval. Use both variance-endpoint branches
in \eqref{ce:variance-endpoint-max}, with the same actual noise
and the lower-allocation moment cost in both. Each branch has
its own whole-interval noise and curvature checks before their
maximum is bounded. The lowest-count square starts at phase
$1/8$, below the contact lower bound $.15468890249$; on the
sixth noise interval, auxiliary phase $13/16$ meets original
count-$12$ phase $-3/16$.

At $30$, contact leaves only $12,13$, with valid outer phase
intervals $[-7/16,1/2]$ and $[-1/2,5/16]$. The negative part
of count $12$ uses a lower-allocation square and current prices;
on the lowest noise interval the same current reference is
interpolated through one extra phase face, $-1/8$.
Its positive part hands off at $1/128$ with width $2/3$ on
that interval. Here $y_{\max}=125/384>0$ and the one-clip
and variance-derivative conditions hold on the whole region.
Higher-noise square prices use width $5/4$ where specified
by the plan, and the shifted-count joins close count $13$.

All these explicit finite plans retain the seven original noise and
four allocation intervals. Their remaining pieces use the current
or square envelopes, and their closed unions cover every retained
contact interval. Their largest upper bounds are those in the table.
\end{proof}
